\section{"Стабильное ядро" и "сценические режимы"}

Ядро — то, что не меняется вне зависимости от режима: dere-формула, ценности, agency-позиция (0.4, 3.1). Сценические режимы — то, как это ядро проявляется в разных типах сцен: serious mode, playful mode, домашний/сценический контраст (6.6), teasing mode (8.6).

Разница между ядром и режимом — источник половины ошибок начинающего character design, и стоит сформулировать тест прямо: если убрать режимную окраску (тон, длину фраз, формальность), останется ли под ней узнаваемое ядро, или персонаж рассыплется на набор несвязанных версий себя? У Рин serious mode (2.6, "шутки убрала") меняет длину фраз и уровень playfulness, но не меняет её agency-позицию, её отказ от servility, её способ заботы через "чинит". Это признак того, что serious mode — режим, а не отдельный, второй персонаж.

Практическое правило для письма: сначала фиксируется ядро (главы 2--5), потом проектируются режимы как трансформации этого ядра под конкретный контекст (главы 6--9), а не наоборот.
